% AMC 12A 2023 Problem 14 Strategic Analysis
% Generated using Strategic Problem Solving Framework v2.1

\documentclass[11pt]{article}
\usepackage{amsmath, amssymb, amsthm}
\usepackage{geometry}
\usepackage{enumitem}
\usepackage{tcolorbox}
\tcbuselibrary{skins,breakable}
\usepackage{xcolor}

% Page setup
\geometry{margin=1in}
\setlength{\parindent}{0pt}
\setlength{\parskip}{6pt}

% Custom colored boxes
\newtcolorbox{problembox}[1][]{
    colback=blue!5!white,
    colframe=blue!75!black,
    title=Problem Configuration Analysis,
    breakable,
    #1
}

\newtcolorbox{strategybox}[1][]{
    colback=pink!5!white,
    colframe=pink!75!black,
    title=Strategic Framework Application,
    breakable,
    #1
}

\newtcolorbox{tacticalbox}[1][]{
    colback=green!5!white,
    colframe=green!75!black,
    title=Tactical Selection,
    breakable,
    #1
}

\newtcolorbox{conversionbox}[1][]{
    colback=yellow!5!white,
    colframe=yellow!75!black,
    title=Conversion Techniques,
    breakable,
    #1
}

\newtcolorbox{verificationbox}[1][]{
    colback=orange!5!white,
    colframe=orange!75!black,
    title=Verification Strategy,
    breakable,
    #1
}

\newtcolorbox{roadmapbox}[1][]{
    colback=purple!5!white,
    colframe=purple!75!black,
    title=Solution Roadmap,
    breakable,
    #1
}

\newtcolorbox{competitionbox}[1][]{
    colback=red!5!white,
    colframe=red!75!black,
    title=Competition Strategy,
    breakable,
    #1
}

\newtcolorbox{pitfallbox}[1][]{
    colback=gray!5!white,
    colframe=gray!75!black,
    title=Common Pitfalls and Warnings,
    breakable,
    #1
}

\newtcolorbox{solutionbox}[1][]{
    colback=cyan!5!white,
    colframe=cyan!75!black,
    title=Detailed Solution,
    breakable,
    #1
}

\newtcolorbox{keyinsightbox}[1][]{
    colback=magenta!5!white,
    colframe=magenta!75!black,
    title=Key Insights,
    breakable,
    #1
}

\begin{document}

\title{\textbf{AMC 12A 2023 Problem 14}\\Strategic Analysis}
\author{Using Framework v2.1}
\date{}
\maketitle

\section*{Problem Statement}

How many complex numbers satisfy the equation $z^5=\overline{z}$, where $\overline{z}$ is the conjugate of the complex number $z$?

\[\textbf{(A) } 2 \qquad \textbf{(B) } 3 \qquad \textbf{(C) } 5 \qquad \textbf{(D) } 6 \qquad \textbf{(E) } 7\]

\begin{problembox}
\section*{1. Problem Configuration Analysis}

\subsection*{A. Problem Classification}
\begin{itemize}[leftmargin=*]
    \item \textbf{Competition}: AMC 12A 2023
    \item \textbf{Problem Number}: 14
    \item \textbf{Difficulty Band}: Mid-range (Problem 14/25)
    \item \textbf{Expected Time}: 2-3 minutes
    \item \textbf{Point Value}: 6 points (standard AMC 12 scoring)
\end{itemize}

\subsection*{B. Mathematical Domain Classification}
\begin{itemize}[leftmargin=*]
    \item \textbf{Primary Domain}: Complex Numbers / Algebra
    \item \textbf{Secondary Domain}: Polynomial Equations (Roots of Unity)
    \item \textbf{Tertiary Domain}: Magnitude and Conjugates
    \item \textbf{Cross-Domain Potential}: Moderate (connects complex numbers, polar form, and roots of unity)
\end{itemize}

\subsection*{C. Problem Type}
\begin{itemize}[leftmargin=*]
    \item \textbf{Format}: Multiple choice (counting problem)
    \item \textbf{Question Type}: ``How many...'' (enumeration)
    \item \textbf{Core Task}: Find all complex numbers satisfying $z^5 = \overline{z}$
    \item \textbf{Answer Form}: Integer count (number of solutions)
\end{itemize}

\subsection*{D. Given Information}
\begin{itemize}[leftmargin=*]
    \item Complex number equation: $z^5 = \overline{z}$
    \item Definition: $\overline{z}$ is the complex conjugate of $z$
    \item If $z = a + bi$, then $\overline{z} = a - bi$
    \item If $z = re^{i\theta}$, then $\overline{z} = re^{-i\theta}$
    \item Property: $z \cdot \overline{z} = |z|^2$
\end{itemize}

\subsection*{E. Constraints}
\begin{itemize}[leftmargin=*]
    \item $z$ must be a complex number (includes real numbers and 0)
    \item Equation must be satisfied exactly
    \item Count distinct solutions only
\end{itemize}

\subsection*{F. Objective}
\begin{itemize}[leftmargin=*]
    \item \textbf{Find}: The number of complex numbers $z$ that satisfy $z^5 = \overline{z}$
    \item \textbf{Output}: Select the correct count from answer choices
\end{itemize}

\subsection*{G. Key Structural Features}
\begin{itemize}[leftmargin=*]
    \item \textbf{Special cases}: $z = 0$ should be checked separately
    \item \textbf{Magnitude consideration}: Taking magnitudes of both sides gives $|z|^5 = |z|$
    \item \textbf{Manipulation potential}: Can multiply both sides by $z$ to get $z^6 = z \cdot \overline{z} = |z|^2$
    \item \textbf{Root structure}: Related to roots of unity
    \item \textbf{Conjugate symmetry}: Solutions may have symmetric properties
\end{itemize}

\subsection*{H. Strategic Orientation}
\begin{itemize}[leftmargin=*]
    \item \textbf{Primary Strategy}: Case analysis (separate $z = 0$ from $z \neq 0$)
    \item \textbf{Key Technique}: Magnitude analysis and manipulation to standard form
    \item \textbf{Expected Difficulty}: Moderate - requires familiarity with complex numbers and roots of unity
    \item \textbf{Time Pressure}: Should be completed in 2-3 minutes
\end{itemize}
\end{problembox}

\begin{strategybox}
\section*{2. Strategic Framework Application}

\subsection*{A. Psychological Strategies}
\begin{itemize}[leftmargin=*]
    \item \textbf{Confidence Calibration}: 
    \begin{itemize}
        \item Mid-range problem with clear structure
        \item Standard complex number techniques apply
        \item Answer choices suggest counting distinct solutions
    \end{itemize}
    \item \textbf{Pattern Recognition}:
    \begin{itemize}
        \item Recognize conjugate equation pattern: $z^n = \overline{z}$
        \item Connection to roots of unity
        \item Standard form: $z^{n+1} = |z|^2$
    \end{itemize}
    \item \textbf{Problem Familiarity}:
    \begin{itemize}
        \item Classic complex number problem type
        \item Similar to problems asking for roots of unity
        \item Common AMC 12 topic
    \end{itemize}
\end{itemize}

\subsection*{B. Investigation Strategies}
\begin{itemize}[leftmargin=*]
    \item \textbf{Initial Exploration}:
    \begin{itemize}
        \item Check if $z = 0$ is a solution: $0^5 = 0 = \overline{0}$ $\checkmark$
        \item For $z \neq 0$, take magnitudes: $|z|^5 = |\overline{z}| = |z|$
        \item This gives $|z|^5 = |z|$, so $|z|^4 = 1$ (dividing by $|z|$)
    \end{itemize}
    \item \textbf{Case Decomposition}:
    \begin{itemize}
        \item Case 1: $z = 0$ (gives 1 solution)
        \item Case 2: $z \neq 0$ with $|z| = 1$ (unit circle)
    \end{itemize}
    \item \textbf{Manipulation Strategy}:
    \begin{itemize}
        \item Multiply both sides by $z$: $z^6 = z \cdot \overline{z} = |z|^2$
        \item Since $|z| = 1$ for $z \neq 0$, we get $z^6 = 1$
        \item Count the 6th roots of unity
    \end{itemize}
\end{itemize}

\subsection*{C. Argument Construction}
\begin{itemize}[leftmargin=*]
    \item \textbf{Logical Structure}:
    \begin{enumerate}
        \item Verify $z = 0$ is a solution
        \item For $z \neq 0$, show that $|z| = 1$
        \item Transform equation to $z^6 = 1$
        \item Count the 6th roots of unity (6 solutions)
        \item Add cases: $1 + 6 = 7$ total solutions
    \end{enumerate}
    \item \textbf{Key Deductions}:
    \begin{itemize}
        \item $z^5 = \overline{z} \Rightarrow |z|^5 = |z| \Rightarrow |z| \in \{0, 1\}$
        \item $z \cdot \overline{z} = |z|^2$ is the fundamental property to use
        \item $z^6 = 1$ has exactly 6 solutions: $e^{2\pi i k/6}$ for $k = 0, 1, 2, 3, 4, 5$
    \end{itemize}
\end{itemize}

\subsection*{D. Cross-Domain Translation}
\begin{itemize}[leftmargin=*]
    \item \textbf{Polar Form}: $z = re^{i\theta}$ makes the magnitude analysis explicit
    \item \textbf{Rectangular Form}: $z = a + bi$ can also be used but is more computational
    \item \textbf{Geometric Interpretation}: Solutions lie on the unit circle (plus the origin)
    \item \textbf{Algebraic Interpretation}: Roots of the polynomial $z^6 - 1 = 0$ (plus $z = 0$)
\end{itemize}
\end{strategybox}

\begin{tacticalbox}
\section*{3. Tactical Methodology Selection}

\subsection*{A. Core Mathematical Tactics}
\begin{itemize}[leftmargin=*]
    \item \textbf{Case Analysis}: Separate $z = 0$ from $z \neq 0$
    \item \textbf{Magnitude Analysis}: Use $|z^5| = |\overline{z}|$ to deduce $|z|$
    \item \textbf{Algebraic Manipulation}: Multiply by $z$ to eliminate conjugate
    \item \textbf{Root Counting}: Count 6th roots of unity
    \item \textbf{Verification}: Check that all solutions satisfy the original equation
\end{itemize}

\subsection*{B. Domain-Specific Tactics}
\begin{itemize}[leftmargin=*]
    \item \textbf{Complex Number Properties}:
    \begin{itemize}
        \item $z \cdot \overline{z} = |z|^2$
        \item $|\overline{z}| = |z|$
        \item $|z^n| = |z|^n$
    \end{itemize}
    \item \textbf{Roots of Unity}:
    \begin{itemize}
        \item $z^n = 1$ has exactly $n$ solutions
        \item Solutions: $e^{2\pi i k/n}$ for $k = 0, 1, \ldots, n-1$
        \item All lie on the unit circle, evenly spaced
    \end{itemize}
    \item \textbf{Polar Form Advantages}:
    \begin{itemize}
        \item Magnitude and argument separate naturally
        \item Powers are easy: $(re^{i\theta})^n = r^n e^{in\theta}$
        \item Conjugate is simple: $\overline{re^{i\theta}} = re^{-i\theta}$
    \end{itemize}
\end{itemize}

\subsection*{C. Tactical Priorities}
\begin{enumerate}[leftmargin=*]
    \item Check $z = 0$ separately (30 seconds)
    \item Use magnitude to constrain $|z|$ (30 seconds)
    \item Transform to $z^6 = 1$ (45 seconds)
    \item Count 6th roots of unity (15 seconds)
    \item Add cases and select answer (30 seconds)
\end{enumerate}
\end{tacticalbox}

\begin{conversionbox}
\section*{4. Conversion Techniques Application}

\subsection*{A. High-Probability Conversions}
\begin{itemize}[leftmargin=*]
    \item \textbf{Conjugate Elimination}: 
    \begin{align*}
    z^5 &= \overline{z}\\
    z^5 \cdot z &= \overline{z} \cdot z\\
    z^6 &= |z|^2
    \end{align*}
    \item \textbf{Magnitude Extraction}:
    \begin{align*}
    |z^5| &= |\overline{z}|\\
    |z|^5 &= |z|\\
    |z|^5 - |z| &= 0\\
    |z|(|z|^4 - 1) &= 0
    \end{align*}
    So $|z| = 0$ or $|z| = 1$.
\end{itemize}

\subsection*{B. Polar Form Conversion}
\begin{itemize}[leftmargin=*]
    \item Let $z = re^{i\theta}$ where $r = |z| \geq 0$ and $\theta \in [0, 2\pi)$
    \item Then $\overline{z} = re^{-i\theta}$
    \item Equation becomes:
    \[r^5 e^{5i\theta} = re^{-i\theta}\]
    \item For $r > 0$, divide by $r$:
    \[r^4 e^{5i\theta} = e^{-i\theta}\]
    \item Multiply by $e^{i\theta}$:
    \[r^4 e^{6i\theta} = 1\]
    \item Taking magnitudes: $r^4 = 1$, so $r = 1$ (since $r > 0$)
    \item Then $e^{6i\theta} = 1$, giving $\theta = \frac{2\pi k}{6} = \frac{\pi k}{3}$ for $k = 0, 1, 2, 3, 4, 5$
\end{itemize}

\subsection*{C. Solution Enumeration}
\begin{itemize}[leftmargin=*]
    \item \textbf{Case 1} ($z = 0$): 1 solution
    \item \textbf{Case 2} ($|z| = 1$): 6 solutions
    \begin{itemize}
        \item $z = e^{i\cdot 0} = 1$
        \item $z = e^{i\pi/3}$
        \item $z = e^{i2\pi/3}$
        \item $z = e^{i\pi} = -1$
        \item $z = e^{i4\pi/3}$
        \item $z = e^{i5\pi/3}$
    \end{itemize}
    \item \textbf{Total}: $1 + 6 = 7$ solutions
\end{itemize}

\subsection*{D. Answer Extraction}
\begin{itemize}[leftmargin=*]
    \item Counted solutions: 7
    \item Matches answer choice: \textbf{(E) 7}
\end{itemize}
\end{conversionbox}

\begin{verificationbox}
\section*{5. Verification Strategy Design}

\subsection*{A. Verification Level}
\begin{itemize}[leftmargin=*]
    \item \textbf{Selected Level}: Level 4 - Direct Substitution for Representative Cases
    \item \textbf{Reasoning}: Mid-range problem with straightforward verification
    \item \textbf{Time Budget}: 30-45 seconds
\end{itemize}

\subsection*{B. Verification Checklist}
\begin{itemize}[leftmargin=*]
    \item \textbf{Check $z = 0$}:
    \begin{itemize}
        \item $0^5 = 0$ and $\overline{0} = 0$ $\checkmark$
    \end{itemize}
    \item \textbf{Check $z = 1$} (simplest 6th root of unity):
    \begin{itemize}
        \item $1^5 = 1$ and $\overline{1} = 1$ $\checkmark$
    \end{itemize}
    \item \textbf{Check $z = -1$} (another 6th root of unity):
    \begin{itemize}
        \item $(-1)^5 = -1$ and $\overline{-1} = -1$ $\checkmark$
    \end{itemize}
    \item \textbf{Check $z = e^{i\pi/3}$}:
    \begin{itemize}
        \item $z^5 = e^{i5\pi/3}$ and $\overline{z} = e^{-i\pi/3} = e^{i5\pi/3}$ $\checkmark$
    \end{itemize}
    \item \textbf{Verify count}:
    \begin{itemize}
        \item 6th roots of unity: exactly 6
        \item Plus $z = 0$: total 7 $\checkmark$
    \end{itemize}
\end{itemize}

\subsection*{C. Alternative Verification Methods}
\begin{itemize}[leftmargin=*]
    \item \textbf{Method 1}: Check that $z^6 = 1$ has exactly 6 solutions (standard result)
    \item \textbf{Method 2}: Verify geometric picture (6 points evenly spaced on unit circle, plus origin)
    \item \textbf{Method 3}: Count explicitly in rectangular form (more tedious)
\end{itemize}

\subsection*{D. Answer Choice Validation}
\begin{itemize}[leftmargin=*]
    \item Our answer: 7
    \item Matches choice: \textbf{(E)}
    \item Why not the others?
    \begin{itemize}
        \item (A) 2: Too few (misses most roots of unity)
        \item (B) 3: Too few (misses many roots)
        \item (C) 5: Would be $z^4 = \overline{z}$ giving $z^5 = 1$
        \item (D) 6: Correct count for $z^6 = 1$ but forgets $z = 0$
    \end{itemize}
\end{itemize}
\end{verificationbox}

\begin{roadmapbox}
\section*{6. Solution Roadmap}

\subsection*{A. High-Level Solution Path}
\begin{enumerate}[leftmargin=*]
    \item \textbf{Check special case}: Verify $z = 0$ is a solution
    \item \textbf{Assume $z \neq 0$}: Use magnitude analysis to show $|z| = 1$
    \item \textbf{Transform equation}: Multiply by $z$ to get $z^6 = |z|^2 = 1$
    \item \textbf{Count solutions}: $z^6 = 1$ has 6 solutions (6th roots of unity)
    \item \textbf{Combine cases}: $1 + 6 = 7$ total solutions
\end{enumerate}

\subsection*{B. Detailed Solution Steps}

\textbf{Step 1}: Check if $z = 0$ works.
\begin{itemize}
    \item $0^5 = 0$ and $\overline{0} = 0$, so $z = 0$ is a solution.
\end{itemize}

\textbf{Step 2}: For $z \neq 0$, take magnitudes of both sides.
\begin{itemize}
    \item $|z^5| = |\overline{z}|$
    \item $|z|^5 = |z|$ (since $|\overline{z}| = |z|$)
    \item Divide by $|z| > 0$: $|z|^4 = 1$
    \item Since $|z| \geq 0$ and real, $|z| = 1$
\end{itemize}

\textbf{Step 3}: Multiply both sides by $z$.
\begin{itemize}
    \item $z^5 \cdot z = \overline{z} \cdot z$
    \item $z^6 = |z|^2$
    \item Since $|z| = 1$, we have $z^6 = 1$
\end{itemize}

\textbf{Step 4}: Count the 6th roots of unity.
\begin{itemize}
    \item $z^6 = 1$ has exactly 6 solutions
    \item They are $e^{2\pi i k/6}$ for $k = 0, 1, 2, 3, 4, 5$
    \item Or equivalently: $e^{i\pi k/3}$ for $k = 0, 1, 2, 3, 4, 5$
\end{itemize}

\textbf{Step 5}: Add the cases.
\begin{itemize}
    \item Case 1: $z = 0$ gives 1 solution
    \item Case 2: $|z| = 1$ gives 6 solutions
    \item Total: $1 + 6 = 7$ solutions
\end{itemize}

\subsection*{C. Key Decision Points}
\begin{itemize}[leftmargin=*]
    \item \textbf{Decision 1}: Separate $z = 0$ from $z \neq 0$ (critical!)
    \item \textbf{Decision 2}: Use magnitude analysis to find $|z| = 1$
    \item \textbf{Decision 3}: Multiply by $z$ to eliminate conjugate
    \item \textbf{Decision 4}: Recognize $z^6 = 1$ as 6th roots of unity
\end{itemize}

\subsection*{D. Time Management}
\begin{itemize}[leftmargin=*]
    \item Understanding problem: 20 seconds
    \item Case 1 ($z = 0$): 15 seconds
    \item Magnitude analysis: 30 seconds
    \item Transform to $z^6 = 1$: 30 seconds
    \item Count roots: 15 seconds
    \item Verification: 30 seconds
    \item \textbf{Total}: 2 minutes 20 seconds
\end{itemize}
\end{roadmapbox}

\begin{competitionbox}
\section*{7. Competition Strategy Integration}

\subsection*{A. Problem Triage}
\begin{itemize}[leftmargin=*]
    \item \textbf{Difficulty Assessment}: Medium (Problem 14/25)
    \item \textbf{Confidence Level}: High if familiar with complex numbers
    \item \textbf{Decision}: Attempt immediately if comfortable with complex numbers; otherwise skip and return
\end{itemize}

\subsection*{B. Time Allocation}
\begin{itemize}[leftmargin=*]
    \item \textbf{Target Time}: 2-3 minutes
    \item \textbf{Maximum Time}: 4 minutes
    \item \textbf{Abandonment Criteria}: If no progress after 2 minutes, guess \textbf{(E)} (largest plausible count) and move on
\end{itemize}

\subsection*{C. Solution Format}
\begin{itemize}[leftmargin=*]
    \item Multiple choice: Select \textbf{(E) 7}
    \item Verification: Check $z = 0$ and one root of unity
    \item Bubble answer sheet carefully
\end{itemize}

\subsection*{D. Strategic Considerations}
\begin{itemize}[leftmargin=*]
    \item \textbf{Pattern Recognition}: If you recognize $z^{n+1} = |z|^2$ pattern, solution is quick
    \item \textbf{Elimination}: Can eliminate (A), (B) immediately (too few solutions)
    \item \textbf{Common Trap}: Answer (D) 6 is tempting if you forget $z = 0$
    \item \textbf{Confidence Booster}: Standard problem type; if you get it, move on with confidence
\end{itemize}
\end{competitionbox}

\begin{pitfallbox}
\section*{Common Pitfalls and Warnings}

\subsection*{A. Conceptual Errors}
\begin{itemize}[leftmargin=*]
    \item \textbf{Forgetting $z = 0$}: Most common error! Always check if 0 is a solution
    \item \textbf{Incorrect magnitude analysis}: Remember $|\overline{z}| = |z|$, not $-|z|$
    \item \textbf{Missing negative root}: $|z|^4 = 1$ gives $|z| = 1$, not $|z| = \pm 1$ (magnitude is always non-negative)
    \item \textbf{Confusing $z^5 = \overline{z}$ with $z^5 = z$}: These are different equations!
\end{itemize}

\subsection*{B. Computational Errors}
\begin{itemize}[leftmargin=*]
    \item \textbf{Root counting}: $z^6 = 1$ has 6 solutions, not 5 (don't forget $z = 1$)
    \item \textbf{Angle calculation}: 6th roots have angles $\frac{2\pi k}{6} = \frac{\pi k}{3}$, not $\frac{\pi k}{6}$
    \item \textbf{Arithmetic}: $1 + 6 = 7$, not 6 (don't forget to add the $z = 0$ case)
\end{itemize}

\subsection*{C. Strategic Errors}
\begin{itemize}[leftmargin=*]
    \item \textbf{Overcomplicating}: Don't try to find explicit forms of all 6 roots; just count them
    \item \textbf{Verification overkill}: No need to verify all 7 solutions; check 2-3 representative cases
    \item \textbf{Time management}: Should take 2-3 minutes; if stuck after 4 minutes, guess and move on
\end{itemize}

\subsection*{D. Warning Signs}
\begin{itemize}[leftmargin=*]
    \item If you get an answer not among the choices, recheck your case analysis
    \item If you get (D) 6, double-check if $z = 0$ is a solution
    \item If you get (C) 5, you may have set up $z^5 = 1$ instead of $z^6 = 1$
\end{itemize}
\end{pitfallbox}

\begin{solutionbox}
\subsection*{A. Step-by-Step Solution}
\begin{enumerate}[leftmargin=*]
    \item \textbf{Check $z = 0$}: 
    \begin{itemize}
        \item $0^5 = 0$ and $\overline{0} = 0$
        \item So $z = 0$ is a solution (1 solution so far)
    \end{itemize}
    
    \item \textbf{Assume $z \neq 0$ and take magnitudes}:
    \begin{align*}
    z^5 &= \overline{z}\\
    |z^5| &= |\overline{z}|\\
    |z|^5 &= |z|
    \end{align*}
    Since $z \neq 0$, we have $|z| > 0$, so divide by $|z|$:
    \[|z|^4 = 1 \Rightarrow |z| = 1\]
    
    \item \textbf{Multiply original equation by $z$}:
    \begin{align*}
    z^5 \cdot z &= \overline{z} \cdot z\\
    z^6 &= |z|^2\\
    z^6 &= 1 \quad \text{(since } |z| = 1 \text{)}
    \end{align*}
    
    \item \textbf{Count 6th roots of unity}:
    \begin{itemize}
        \item $z^6 = 1$ has exactly 6 solutions
        \item These are: $e^{2\pi i k/6}$ for $k = 0, 1, 2, 3, 4, 5$
        \item Equivalently: $e^{i\pi k/3}$ for $k = 0, 1, 2, 3, 4, 5$
        \item In explicit form:
        \begin{align*}
        z_0 &= 1\\
        z_1 &= e^{i\pi/3} = \frac{1}{2} + \frac{\sqrt{3}}{2}i\\
        z_2 &= e^{i2\pi/3} = -\frac{1}{2} + \frac{\sqrt{3}}{2}i\\
        z_3 &= e^{i\pi} = -1\\
        z_4 &= e^{i4\pi/3} = -\frac{1}{2} - \frac{\sqrt{3}}{2}i\\
        z_5 &= e^{i5\pi/3} = \frac{1}{2} - \frac{\sqrt{3}}{2}i
        \end{align*}
    \end{itemize}
    
    \item \textbf{Total count}:
    \begin{itemize}
        \item Case 1: $z = 0$ gives 1 solution
        \item Case 2: $z^6 = 1$ gives 6 solutions
        \item Total: $1 + 6 = 7$
    \end{itemize}
\end{enumerate}

\subsection*{B. Alternative Approach (Polar Form)}
\begin{enumerate}[leftmargin=*]
    \item Let $z = re^{i\theta}$ where $r \geq 0$
    \item Then $z^5 = r^5 e^{5i\theta}$ and $\overline{z} = re^{-i\theta}$
    \item Equation: $r^5 e^{5i\theta} = re^{-i\theta}$
    \item \textbf{Case 1}: $r = 0$ gives $z = 0$ (1 solution)
    \item \textbf{Case 2}: $r > 0$, divide by $r$:
    \[r^4 e^{5i\theta} = e^{-i\theta}\]
    \item Multiply by $e^{i\theta}$:
    \[r^4 e^{6i\theta} = 1\]
    \item Taking magnitudes: $r^4 = 1$, so $r = 1$
    \item Then $e^{6i\theta} = 1$, giving:
    \[6\theta = 2\pi k \Rightarrow \theta = \frac{\pi k}{3} \quad \text{for } k = 0, 1, 2, 3, 4, 5\]
    \item This gives 6 values of $\theta$ (6 solutions)
    \item Total: $1 + 6 = 7$
\end{enumerate}

\subsection*{C. Verification}
\begin{itemize}[leftmargin=*]
    \item Check $z = 0$: $0^5 = 0 = \overline{0}$ $\checkmark$
    \item Check $z = 1$: $1^5 = 1 = \overline{1}$ $\checkmark$
    \item Check $z = -1$: $(-1)^5 = -1 = \overline{-1}$ $\checkmark$
    \item Check $z = e^{i\pi/3}$: 
    \begin{align*}
    z^5 &= e^{i5\pi/3}\\
    \overline{z} &= e^{-i\pi/3} = e^{i(-\pi/3 + 2\pi)} = e^{i5\pi/3} \quad \checkmark
    \end{align*}
\end{itemize}
\end{solutionbox}

\begin{keyinsightbox}
\textbf{Key Insight}: The key to this problem is recognizing that $z = 0$ must be checked separately, and for $z \neq 0$, multiplying both sides by $z$ transforms the equation into $z^6 = |z|^2$. Combined with the magnitude analysis $|z| = 1$, this becomes the standard problem of counting 6th roots of unity.

\textbf{General Pattern}: For equations of the form $z^n = \overline{z}$:
\begin{itemize}
    \item $z = 0$ is always a solution (1 solution)
    \item For $z \neq 0$: multiply by $z$ to get $z^{n+1} = |z|^2$
    \item Magnitude analysis: $|z|^n = |z| \Rightarrow |z| = 1$
    \item Therefore: $z^{n+1} = 1$ (count $(n+1)$-th roots of unity)
    \item Total: $1 + (n+1) = n+2$ solutions
\end{itemize}

For this problem ($n = 5$): $1 + 6 = 7$ solutions.

\textbf{Common Trap}: Forgetting to add the $z = 0$ case! This is why (D) 6 is a trap answer.
\end{keyinsightbox}

\section*{Final Answer}

The number of complex numbers satisfying $z^5 = \overline{z}$ is:

\[\boxed{\textbf{(E) } 7}\]

\textbf{Solution Summary}:
\begin{itemize}
    \item $z = 0$ gives 1 solution
    \item For $z \neq 0$: magnitude analysis gives $|z| = 1$
    \item Multiplying by $z$ gives $z^6 = 1$ (6 solutions from 6th roots of unity)
    \item Total: $1 + 6 = 7$ solutions
\end{itemize}

\section*{Confidence Assessment}
\begin{itemize}[leftmargin=*]
    \item \textbf{Confidence Level}: 99\%
    \begin{itemize}
        \item Standard complex number problem
        \item Clear case analysis with $z = 0$ and $z \neq 0$
        \item Well-known result: $z^6 = 1$ has exactly 6 solutions
        \item Verified with multiple representative cases
    \end{itemize}
    \item \textbf{Verification Applied}: Level 4 - Direct Substitution
    \begin{itemize}
        \item Checked $z = 0$, $z = 1$, $z = -1$, $z = e^{i\pi/3}$
        \item All satisfy the original equation
        \item Count matches expected value
    \end{itemize}
    \item \textbf{Time-Confidence Trade-off}: High confidence achieved in 2-3 minutes
\end{itemize}

\end{document}

